\section{Four Lessons from Two Years of Use}
\label{sec:lessons}

We designed Clay~v0, released in early 2024, as a masked autoencoder for end-to-end fine-tuning, not as a standalone embedding model. In practice, embeddings have become one of the primary use cases for the Clay model. Chip embeddings are small: a $1{,}024$-D float16 vector takes $2\,\text{KB}$, against $1.3\,\text{MB}$ for the 10-band int16 $256\times 256$ image chip it summarizes (${\sim}640\times$ smaller), so screening large areas or long time series can run on stored vectors instead of imagery ~\cite{gomes2024nec,vinge2025neucobench}. We do not compare chip-level and pixel-level accuracy per unit cost here. Public GeoParquet catalogs on Source Cooperative (\href{https://source.coop/clay}{\texttt{https://\allowbreak source.\allowbreak coop/\allowbreak clay}})~\cite{clay2026sourcecoop} now serve over $15.2\,\text{B}$ global Sentinel-2 Clay embeddings across $2.8\,\text{M}$ scenes (2017--2026, \texttt{clay/\allowbreak lgnd-\allowbreak embeddings}) and CONUS-wide 2010--2021 NAIP~v1.5 embeddings (\texttt{clay/\allowbreak clay-\allowbreak v1-\allowbreak 5-\allowbreak naip-\allowbreak 2}). The NAIP catalog was chipped with \texttt{stacchip}.

These lessons draw in part on work by the authors applying Earth embeddings on over $20$ projects in which we partnered with nonprofits, companies and a development bank to turn open satellite data and AI into working answers to their questions, taking each from a first conversation through discovery and prototype toward a build they could run themselves~\cite{grimshaw2026soy,chen2026aquaculture}. Partners include Amazon Sustainability, Aquarry, the Asian Development Bank, Carbon Mapper, CTrees, the Jane Goodall Institute, The Nature Conservancy, Seafood Watch, Trase, and Xylo Earth.

\paragraph{1. Frozen backbones with spatial decoders come within $0.5$ IoU of full fine-tuning on flood segmentation; neither separates from a from-scratch U-Net.}
\emph{Takeaway:} keep the backbone frozen and train a small spatial decoder when boundaries matter; a per-token linear head is too coarse at $10\,\text{m}$. \label{sec:lesson_finetune}
On 10-band Sentinel-2 flood segmentation (\texttt{Sen1Floods11}~\cite{bonafilia2020sen1floods11}), a $1\times 1$ linear head on frozen $80\,\text{m}$ tokens blurs $10\,\text{m}$ shorelines ($74.8$ water IoU vs.\ $83.1$ for a $31\text{M}$-parameter from-scratch U-Net). A lightweight UPerNet decoder~\cite{xiao2018upernet} ($8.9\text{M}$ params) on the \textbf{frozen} Clay encoder raises water IoU to $83.3$ with $1/36$ of the trainable parameters. Full fine-tuning ($320.3\text{M}$ parameters) reaches $83.8$ (difference $+0.5$, paired $95\%$ interval $[-0.1, +1.5]$). Consistent with PANGAEA~\cite{marsocci2026pangaea}, foundation models barely edge out a supervised U-Net at $100\%$ labels. With few labels, others report clear gains: on 5-shot flood mapping~\cite{kaushik2026flood}, and on $11$-class tropical land-cover segmentation, where frozen Clay~v1 with a decoder gains $+7.5\text{--}11.4$ mIoU over U-Net++ and Prithvi~\cite{wiratama2025indonesia}.

\paragraph{2. \texttt{CLS} summarizes the scene, patch tokens keep the layout, and sensors do not share a common semantic space.}
\emph{Takeaway:} use \texttt{CLS} for scene-level search and classification, patch tokens when position within the chip matters, and do not expect a probe trained on one sensor to transfer to another. \label{sec:lesson_tokens}
Pooling trade-offs~\cite{corley2026pooling} are sharper in Clay~v1.5 because Stage~2 distilled \texttt{CLS} toward an RGB teacher while supervising the patch tokens $\bm{Z}_{\mathrm{patch}}$ through multi-band pixel reconstruction (Figure~\ref{fig:arch}). The patch tokens retain local radiometry and sub-chip layout. They recover within-chip spectral indices ($R^2 \approx 0.90$), solve $2\times 2$ quadrant localization that \texttt{CLS} mostly misses ($99\%$ vs.\ $42\%$; chance $25\%$), and beat broadcast \texttt{CLS} by $+6\text{--}16$ macro-$F_1$ points on ESA WorldCover across Sentinel-1/2, Landsat, and NAIP. Central $160\,\text{m}$ patch tokens likewise beat full-chip mean pooling for Clay and OlmoEarth (Table~\ref{tab:master_bench}, Panel~B).

\texttt{CLS} summarizes the scene. In USPVDB solar search~\cite{fujita2023uspvdb}, mean pooling dilutes compact arrays into background scrubland ($\text{AP} = 0.02$). \texttt{CLS} on $256\times 256$ chips preserves detection ($\text{AP} = 0.27$), as long as chips are not tightly cropped~\cite{corley2024resize}.

We identified three further challenges. (i)~\textbf{Shared weights do not align sensors into a common space.} In our own test across $6{,}592$ co-located CONUS sites on $320\,\text{m}$ footprints ($480\,\text{m}$ for Landsat), zero-shot transfer of a Sentinel-2 \texttt{CLS} probe to NAIP, Landsat, or Sentinel-1 lowers macro-$F_1$ by $41\text{--}57$ points. A $1{,}000$-pair Procrustes rotation recovers $0.19\text{--}0.40$ macro-$F_1$, against $0.56\text{--}0.83$ native. Chip size contributes: Sentinel-2 to NAIP drops $22$ points at $2.56\,\text{km}$ and $59$ at $320\,\text{m}$ (see also~\cite{tamazyan2025geocrossbench,marsocci2025xstars}). (ii)~\textbf{Clay treats time as frozen single-date snapshots}, with only a $[\sin,\cos]$ timestamp prior, so phenology and change semantics must be computed post hoc as differences or derivatives across embedding time series. (iii)~\textbf{Clay lacks native vision--language alignment}, so open-vocabulary retrieval and descriptive tasks require fitting a post-hoc text bridge or supervised probe.

\paragraph{3. Do not slice the first dimensions.}
\emph{Takeaway:} keep all $1{,}024$ dimensions; the first $k$ carry no more information than $k$ random ones. \label{sec:lesson_mrl}
In Stage~1, $7$ nested Matryoshka heads (\texttt{dolls=[16..1024]}) forced the first $16$ coordinates of \texttt{CLS} to carry $8.0\times$ higher variance than trailing dimensions. Stage~2 replaced those heads with a training-time \texttt{nn.Linear(1024, 1024)} projection $\bm{P}$ (\S\ref{sec:lineage}; Figure~\ref{fig:arch}). This left Stage~1's $8.0\times$ variance spike in \texttt{dims[0:16]} but rotated the discriminative features across all $1{,}024$ dimensions. On spatially blocked NLCD regression over $223$ Texas scenes~\cite{dewitz2023nlcd}, naive prefix slicing (\texttt{dims[0:64]}) drops $R^2$ from $0.83$ (full $1{,}024\text{-D}$) to $0.62$, the same as $64$ random coordinates. \texttt{CLS} also linearly encodes its own per-sample pretraining reconstruction loss, so a Ridge probe recovers it without running the decoder (\textbf{ELLE}~\cite{nuno2026elle}).

\paragraph{4. Regularized parametric probes (LSVM, logistic/Ridge, 2-layer MLP) beat centroids, $k$-NN, trees, and 2D UMAP.}
\emph{Takeaway:} start with a linear SVM or logistic regression on the embeddings before trying nearest neighbors, trees, or 2D projections. \label{sec:lesson_probes}
On spatially blocked EuroSAT~\cite{helber2019eurosat} ($50\text{--}200$ labels/class), regularized parametric probes (squared-hinge LSVM~\cite{cortes1995svm,fan2008liblinear}, $\ell_2$ logistic regression, and $2$-layer MLPs) all reach $91\text{--}94\%$ macro-$F_1$. They beat $\ell_2$-normalized RandomForest and HistGBDT by $3\text{--}5$ points, cosine $k$-NN by $5$ points, class centroids by $8$ points, and train-only 2D UMAP~\cite{mcinnes2018umap} $+ k$-NN by $8\text{--}13$ points. On continuous NLCD canopy regression, 2D UMAP also drops $R^2$ from $0.85\text{--}0.87$ to $0.62$. In imbalanced $1{:}10$ concept search under $30\%$ label noise, $2$-layer MLPs and LSVMs ($\sim 0.90\text{ AP}$) also beat nearest-neighbor Max-Sim~\cite{johnson2021faiss} ($0.37\text{ AP}$) by learning negative-class boundaries.
